% !TeX root = ../main.tex
%TC:envir tabular [] word
%TC:envir table [] word
\chapter{Introduction}
The trajectory of processor development has shifted fundamentally in recent decades. Due to physical limitations, the era of rapid single-core performance gains effectively ended. In response, manufacturers pivoted towards multi-core architectures. %cite
Software performance gains are no longer a "free lunch" provided by hardware improvements\cite{free_lunch}. To utilize modern processors to their greatest extent, writing concurrent code is essential.

However, moving from a sequential to a concurrent execution model introduces significant cognitive overhead. Managing multiple threads requires a departure from the intuitive, linear flow of traditional programming. This complexity leads to various hard-to-catch concurrency errors, which are notoriously subtle and difficult to reproduce.

This project designs and implements a language named \textit{Coherence}. \textit{Coherence} presents a novel combination of ideas from Pony, Kappa and Gudka's Lockguard (described in section \ref{sec:related_work}) to create a language suited for writing safe and expressive concurrent programs.

\section{Motivation} \label{sec:intro_motivation}
Concurrent code carries the risk of two types of errors: \textit{data races} and \textit{deadlocks}.

Data races occur when two or more threads attempt to access the same memory location simultaneously, where at least one of those accesses is a write. Without synchronization, the final state of the memory becomes dependent on the unpredictable timing of the processor's scheduler. 

The most common defence against data races are \textit{locks}. Locks are synchronization primitives that grant a thread exclusive access to a resource. By locking a piece of data, a thread ensures that no other thread can modify it until the lock is released. However, without a robust locking paradigm, programs can encounter \textit{deadlocks}. Deadlocks occur when two or more threads enter a circular dependency, where each thread holds a lock that another thread needs to proceed. As no thread can move forward or release its currently held lock, the system reaches a state of permanent stagnation.

One programming model that circumvents these problems is the \textit{actor model}. In this paradigm, the \textit{actor} is the fundamental unit of concurrency. Unlike traditional threads, actors do not share memory with each other. Instead, each actor manages a private state and communicates with other actors exclusively through asynchronous message passing. Each actor possesses an internal \textit{mailbox}, which is a FIFO queue that stores incoming messages. These messages are processed sequentially. Because there is no interleaving of execution within an individual actor and no memory is shared between them, the model avoids both data races and deadlocks.

However, this pure actor model has several real costs. Copying data between actors can be expensive, and coordination patterns that are naturally expressed with locks become awkward or incur extra latency. In short, forbidding shared memory and locks entirely simplifies correctness but can hinder performance and expressiveness.

This project aims to demonstrate that safety does not need to come at the cost of performance and expressiveness. \textit{Coherence} is designed to retain the safety and simplicity of the actor model while safely reintroducing shared memory and locks.

\section{Related Work} \label{sec:related_work}

This section briefly introduces three prior works that \textit{Coherence} synthesises. A more detailed technical treatment of each is in Chapter 2.

\subsection{The Pony Programming Language}
The \textit{Pony} programming language\cite{pony_language} safely adds shared-memory into the actor model. The core innovation of Pony is its system of \textit{reference capabilities}. In Pony, every \textit{reference}\footnote{The terms \textit{pointer} and \textit{reference} will be used interchangeably throughout this document.} type has an associated capability that restricts how the reference can be used. The programmer assigns these capabilities to references they use, and the Pony compiler uses these capabilities to statically forbid programs that may lead to data-races. 

A major shortcoming of Pony is its lack of locks. Many coordination patterns (for example, multiple actors accessing a shared data structure) remain awkward without locks (see \Cref{sec:pony_expressive_power}). To safely integrate locks in an actor-based system, \textit{Coherence} draws inspiration from the following two works.

\subsection{Castegren's Kappa}
Castegren introduced a new reference capability called \texttt{locked} in his work on Kappa\cite{kappa_system}. A reference guarded by this capability is equipped with an implicit lock and every access to the reference automatically acquires and releases the lock, preventing data-races.

\subsection{Gudka's Lock Inference for Java}
In his work on Lockguard\cite{gudka_thesis}, Gudka extended Java with \textit{atomic sections} (code blocks that must execute atomically) and designed a system to automatically infer any locks required for atomic execution, shifting the responsibility of taking locks safely from the programmer to the compiler. 

\section{Project Objectives} \label{sec:project_objectives}
The goal of this project is to develop \textit{Coherence}, a language for writing safe and expressive concurrent programs.

\begin{enumerate}
    \item \textit{Coherence}'s type system should rule out programs containing common concurrency hazards, specifically data-races and deadlocks.
    \item Beyond the type system, \textit{Coherence} should be a complete language that allows programmers to actually write and run programs.
    \item \textit{Coherence} must allow expressive concurrent programming. 
\end{enumerate}